
% Robby Cochran - ATS-oriented LaTeX resume
% Build: xelatex robby_cochran_resume.tex
\documentclass[9pt,letterpaper]{extarticle}

\usepackage[letterpaper,margin=0.44in]{geometry}
\usepackage{enumitem}
\usepackage[hidelinks]{hyperref}
\usepackage{titlesec}
\usepackage{tabularx}
\usepackage{array}
\usepackage[T1]{fontenc}
\usepackage{lmodern}
\usepackage{microtype}
\usepackage{xcolor}
\usepackage{ragged2e}

\pagestyle{empty}
\setlength{\parindent}{0pt}
\setlength{\parskip}{0pt}
\setlength{\tabcolsep}{0pt}
\renewcommand{\baselinestretch}{0.93}
\setlist[itemize]{leftmargin=1.05em,itemsep=0.4pt,topsep=1.2pt,parsep=0pt,partopsep=0pt}
\titleformat{\section}{\large\bfseries\uppercase}{}{0pt}{}[\titlerule]
\titlespacing*{\section}{0pt}{3pt}{2pt}

\newcommand{\role}[4]{%
  \textbf{#1} \hfill {\small #2}\\[-1pt]
  {\small #3} \hfill {\small #4}\vspace{1pt}
}
\newcommand{\link}[2]{\href{#1}{#2}}

\begin{document}

{\LARGE \textbf{Robby Cochran}}\\[-1pt]
{\small Technical engineering leader - Security systems, Kubernetes, AI infrastructure, developer platforms}\\[-1pt]
{\small San Francisco, CA \;|\; \link{mailto:kudzoo@gmail.com}{kudzoo@gmail.com} \;|\; \link{https://www.robbycochran.com}{robbycochran.com} \;|\; \link{https://github.com/robbycochran}{github.com/robbycochran} \;|\; \link{https://www.linkedin.com/in/robertacochran}{linkedin.com/in/robertacochran}}

\section*{Profile}
Technical engineering leader with a CS PhD, Staff-level systems/security depth, and people-management experience across startup and enterprise environments. Builds and leads work in technically hard domains: Kubernetes security, eBPF observability, developer platforms, software verification, and AI-agent infrastructure. Open to engineering leadership, Staff+/Principal, forward-deployed engineering, and founder-path roles where product judgment, customer proximity, and deep technical execution all matter.

\section*{Selected Impact}
\begin{itemize}
  \item Built and led runtime-security product work for StackRox/RHACS across behavioral detection, process, network, and file activity, baselines, Kubernetes-native policy, and analyst/developer workflows.
  \item Helped carry StackRox from early startup product through Red Hat acquisition and into RHACS/OpenShift enterprise scale; product context included \$0 to \$50M+ revenue and 100+ releases.
  \item Led ACS agentic foundations: internal hackathon, CI tooling, support triage bots, skills strategy, secure harness primitives, and developer paths from local sandbox to OpenShift execution.
  \item Worked with large customers and design partners to translate deployment constraints into roadmap decisions around scale, compliance, performance, platform support, IBM Power/Z, and false-positive reduction.
  \item Current focus: runtime contracts, evaluation harnesses, observability, and policy controls for AI agent and developer infrastructure.
\end{itemize}

\section*{Experience}
\role{IBM / Red Hat}{2021 - Present}{Staff Engineer / Engineering Manager, Red Hat Advanced Cluster Security (RHACS)}{San Francisco, CA / Remote}
\begin{itemize}
  \item Led RHACS runtime-security roadmap spanning process activity, network security, file activity, behavioral baselines, runtime anomaly detection, Kubernetes-native policy workflows, and developer/analyst experience.
  \item Drove eBPF-based runtime visibility design for OpenShift/Kubernetes environments, including collection semantics, policy exceptions, performance constraints, platform support, and low false-positive detection goals.
  \item Built and guided secure AI-agent infrastructure through \link{https://github.com/stackrox/harness-openshell}{harness-openshell}: declarative workflows, provider/credential boundaries, inference routing, sandbox execution, and repeatable CI validation.
  \item Managed distributed engineering teams while staying close to architecture and implementation; promoted L4 and L5 engineers through scope design, feedback, and career planning.
  \item Shaped OpenShift-first RHACS strategy around console integration, RBAC alignment, GitOps/policy-as-code, runtime observability, and AI-assisted workflows. Public talk: \link{https://speakerdeck.com/redhatlivestreaming/stackrox-community-office-hours-e2-ebpf-101-implementing-security-and-monitoring-kubernetes}{eBPF 101}.
\end{itemize}

\role{StackRox, Inc. (acquired by Red Hat)}{2017 - 2021}{Staff Software Engineer / Senior Software Engineer / Member of Technical Staff}{Mountain View, CA}
\begin{itemize}
  \item Joined an early-stage Kubernetes security startup and built runtime detection and enforcement features for containerized workloads and Kubernetes clusters.
  \item Advanced core runtime-security capabilities including baseline-driven detection, process and network visibility, and connections between observed runtime behavior and security policy.
  \item Worked across product, engineering, and customer-driven security cases with large design partners, helping mature the platform through acquisition by Red Hat.
  \item Patent: \link{https://patents.justia.com/patent/11811804}{US 11,811,804}, ``System and method for detecting process anomalies in a distributed computation system utilizing containers.'' Talk: \link{https://www.youtube.com/watch?v=LFqWl7XB5Gc}{BSidesSF 2018 - Listen to your Engine}.
\end{itemize}

\role{University of North Carolina at Chapel Hill}{2008 - 2016}{PhD Researcher, Computer Security and Software Verification}{Chapel Hill, NC}
\begin{itemize}
  \item Developed symbolic client verification: server-side techniques for checking whether observed client network traffic could have been produced by sanctioned software.
  \item Built symbolic-execution systems for online games, cryptographic protocols, and distributed applications; refactored KLEE internals for parallelism and performance research.
  \item Publications: NDSS 2010 Best Paper / TISSEC 2011 on online-game client verification; NDSS 2013 on distributed-application client behavior; NSDI 2017 on known cryptographic clients.
\end{itemize}

\role{Microsoft Research - Research in Software Engineering (RiSE)}{Summer 2013}{Research Intern}{Redmond, WA}
\begin{itemize}
  \item Program synthesis via crowdsourcing: combined symbolic automata, genetic programming, and online workers; publication: POPL 2015 ``Program Boosting''; patent: \link{https://patents.google.com/patent/US9753696B2}{US 9,753,696}.
\end{itemize}

\role{Intel Corporation - Visual Computing Group}{2006 - 2007}{Software Engineer}{Hillsboro, OR}
\begin{itemize}
  \item Built graphics-hardware analysis tooling to interpret, record, and replay Direct3D/OpenGL API calls for performance analysis of prototype visual-computing hardware.
\end{itemize}

\section*{Education}
\textbf{PhD, Computer Science}, University of North Carolina at Chapel Hill, 2016 - Advisor: Michael K. Reiter.\quad
\textbf{BS, Computer Science}, Clemson University, Magna Cum Laude, 2006.

\section*{Technical Domains}
{\small Runtime security; Kubernetes; OpenShift; RHACS; eBPF observability; behavioral baselines; policy-as-code; distributed systems; symbolic execution; software verification; formal methods; program synthesis; AI-agent safety; LLM gateways; MCP; Go; C/C++; Python; Rust; Bash; Docker; LLVM; SMT/SAT solvers.}

\end{document}
